\documentclass[12pt,a4paper]{article}
\usepackage[utf8]{inputenc}
\usepackage[vietnamese]{babel}
\usepackage[left=2cm,right=2cm,top=2cm,bottom=2cm]{geometry}
\usepackage{mathptmx}
\usepackage{enumitem}

\pagestyle{plain}
\linespread{1.15}

\begin{document}

% --- HEADER QUỐC HIỆU ---
\noindent
\begin{minipage}[t]{0.45\textwidth}
    \begin{center}
        \small
        BỘ CÔNG AN\\
        \textbf{\underline{VIỆN KHOA HỌC HÌNH SỰ}}\\[0.1cm]
        Số: 09/KL-KTHS\\
        \textit{V/v kết luận giám định tài liệu}
    \end{center}
\end{minipage}
\hfill
\begin{minipage}[t]{0.52\textwidth}
    \begin{center}
        \small
        \textbf{CỘNG HÒA XÃ HỘI CHỦ NGHĨA VIỆT NAM}\\
        \textbf{\underline{Độc lập - Tự do - Hạnh phúc}}\\[0.1cm]
        \textit{Hà Nội, ngày 25 tháng 07 năm 2026}
    \end{center}
\end{minipage}

\vspace{0.6cm}

\begin{center}
    \textbf{\large KẾT LUẬN GIÁM ĐỊNH KỸ THUẬT HÌNH SỰ}\\
    \textit{(V/v Giám định tính thật giả của tờ di chúc ký hiệu f2-1)}
\end{center}

\vspace{0.3cm}

\noindent Viện Khoa học Hình sự Bộ Công an kết luận về mẫu giám định di chúc gửi kèm Quyết định trưng cầu số 14/QĐ-CSHS như sau:

\section*{I. NỘI DUNG YÊU CẦU GIÁM ĐỊNH}
Giám định dấu vết tẩy xóa hóa chất, tính đồng nhất của nét chữ, chất liệu mực và thời điểm viết đoạn văn bản chèn đè *"toàn quyền sở hữu và hưởng thụ tiền đền bù"* trong tờ di chúc lập ngày 15/04/2018.

\section*{II. KẾT QUẢ GIÁM ĐỊNH}
\begin{enumerate}[label=\arabic*., leftmargin=1.5cm]
    \item \textbf{Phân tích bề mặt giấy dưới kính phổ quang hồng ngoại/UV:} Phát hiện tổn thương bề mặt xơ sợi giấy tại khu vực dòng tên Trần Ngọc Mai do tác động của dung dịch tẩy xóa hóa chất chuyên dụng (Bleaching agent).
    \item \textbf{Phân tích phổ quang quang học & Tuổi mực:} Đoạn văn bản chèn đè được viết bằng loại mực bi hóa dầu nhãn hiệu hiện đại sản xuất năm 2024, viết đè trực tiếp lên lớp mực bút máy gốc của cụ Nguyễn Văn Thành lập năm 2018.
\end{enumerate}

\section*{III. KẾT LUẬN GIÁM ĐỊNH}
Di chúc lập ngày 15/04/2018 đã bị **TẨY XÓA HÓA CHẤT & VIẾT CHÈN ĐÈ GIẢ MẠO** vào khoảng đầu năm 2024 nhằm xóa tên đồng thừa kế Trần Ngọc Mai để thâu tóm toàn bộ tiền bồi thường đền bù nhà đất.

\vspace{0.8cm}

\noindent
\begin{minipage}[t]{0.45\textwidth}
    \small
    \textbf{\textit{Nơi nhận:}}\\
    - Cơ quan CSĐT Công an TP;\\
    - Lưu: VT, CSHS.
\end{minipage}
\hfill
\begin{minipage}[t]{0.5\textwidth}
    \begin{center}
        \small
        \textbf{VIỆN TRƯỞNG VIỆN KHHS}\\
        \textit{(Ký và đóng dấu)}\\[1.5cm]
        \textbf{Thiếu tướng Trần Quốc Tuấn}
    \end{center}
\end{minipage}

\end{document}
